\section{应用与联系}\label{sec:applications}

\subsection{经典丢番图问题的解决}

纲领最著名的应用是\textbf{费马大定理}：怀尔斯的证明路线（弗雷曲线 $\to$ 里贝特定理 $\to$ 模性提升）正是“伽罗瓦表示嵌入自守世界”的纲领式操作，本文库费马大定理综述已有详述。同类技术随后攻克：

\begin{itemize}
    \item \textbf{Sato--Tate 猜想}：椭圆曲线 Frobenius 角的分布律（2008）；
    \item \textbf{Serre 模性猜想}：奇二维表示的模性（Khare--Wintenberger，2008\cite{khare2009}）；
    \item \textbf{完全平方问题}与同余数等经典问题的进展（借助模形式的系数估计）。
\end{itemize}

\subsection{解析数论的直接红利}

函子性的每一次推进都直接转化为 $L$ 函数与素数分布的定量改进：

\begin{itemize}
    \item \textbf{Ramanujan--Petersson}：对称幂提升\cite{kimshahidi2002}给出 Hecke 特征值的最优级界，是当前进入“过临界带”的主流途径；
    \item \textbf{素数分布}：自守 $L$ 函数的零点自由区域（零密度估计）依赖函子性的单调性结果，间接改进算术级数中素数分布的常数；
    \item \textbf{中心 $L$ 值}：Gan--Gross--Prasad 理论\cite{ggp2012}把分支律与中心临界 $L$ 值的非零性联系起来，是 BSD 猜想与 Gross--Zagier 型公式在高维群上的推广框架。
\end{itemize}

\subsection{表示论与代数几何}

纲领重塑了表示论的版图：Harish-Chandra 模块、Hermite 局部常数、$p$-adic 表示论（$p$-adic Hodge 理论与局部朗兰兹的结合）都以纲领问题为驱动。几何一侧，Shtuka、Hitchin 丛与 Shimura 簇的几何成为代数几何的活跃分支；Ng\^o 对基本引理的证明\cite{ngo2010}展示了周变换在算术问题上的威力，Fargues--Scholze 的几何化\cite{fargues2021}则把完美胚（perfectoid）技术与朗兰兹对应焊接在一起。

\subsection{数学物理}

几何朗兰兹纲领在物理中有一张精确的对照表：Kapustin--Witten 指出，$\mathcal{N}=4$ 超对称杨--米尔斯理论的 \textbf{S-对偶}对应 $\operatorname{Bun}_G$ 的 D-模块与 $^LG$ 局部系统的几何朗兰兹对应；镜像对称中的同调镜像对称亦可视为几何朗兰兹在 Calabi--Yau 情形的表亲。2024年几何朗兰兹的证明\cite{gaitsgoryraskin2024}为这批物理预期提供了第一个严格实例，并推动“范畴化朗兰兹”成为表示论与量子场论共同的通用语言。
